\chapter{Algorithmique et programmation}
\label{corriges:19}
\section*{Exercice 19.1 -- Trace et condition}
Pour $x=-2$, $y=-5$ puis $5$. Pour $x=0$, $y=-1$ puis $1$. Pour $x=3$, $y=5$ puis $25$. Le cas limite est $y=0$, obtenu pour $x=1/2$ ; la condition $y>0$ est fausse, donc la sortie reste $0$. Les valeurs doivent être suivies dans l'ordre des affectations.

\section*{Exercice 19.2 -- Boucle et correction}
Les sommes successives valent $1,3,6,10$ ; à la fin $s=10$ et $i=5$. En remettant $s$ à zéro à chaque tour, il ne resterait que le dernier terme : $4$.
\begin{lstlisting}
mettre n à 0
mettre p à 1
répéter jusqu'à (p >= 100)
    mettre p à (2 * p)
    ajouter 1 à n
fin répéter
\end{lstlisting}
À la fin, $n=7$, $p=128$ ; $2^6=64<100$ garantit qu'il s'agit du premier entier.

\section*{Exercice 19.3 -- Motif et boucle imbriquée}
\begin{lstlisting}
définir triangle (côté)
    répéter 3 fois
        avancer de côté
        tourner à gauche de 120 degrés
    fin répéter

répéter 6 fois
    triangle (60)
    tourner à gauche de 60 degrés
fin répéter
\end{lstlisting}
Initialiser stylo, position et direction. Le bloc revient à son départ après un tour de $360\degre$ ; la boucle extérieure répartit six secteurs de $60\degre$.

\section*{Exercice 19.4 -- Simulation et scripts parallèles}
Initialiser un compteur à $0$ ; répéter $1000$ fois deux tirages entiers séparés entre $1$ et $6$ ; si leur somme vaut $7$, augmenter le compteur ; afficher compteur/$1000$. La valeur théorique à comparer est $1/6$.
Pour le jeu, confier la mise à jour du score à un seul lutin. Un verrou de contact ou l'attente de la séparation empêche plusieurs incréments pendant un même contact. L'autre lutin peut gérer les déplacements dans un script déclenché par le même message de départ.
